\begingroup\singlespacing\fontsize{9.2}{11.32}\selectfont
\setlength{\tabcolsep}{3pt}
\begin{longtable}{@{}>{\raggedright\arraybackslash}p{0.31\linewidth}>{\raggedright\arraybackslash}p{0.12\linewidth}>{\raggedright\arraybackslash}p{0.22\linewidth}>{\raggedright\arraybackslash}p{0.23\linewidth}@{}}
\caption{探索性外层折归一化预测误差}\label{tab:S7}\\
\toprule
程序 & 结局 & 均值 & 折间标准差 \\
\midrule
\endfirsthead
\multicolumn{4}{l}{表 \thetable{} （续）}\\
\toprule
程序 & 结局 & 均值 & 折间标准差 \\
\midrule
\endhead
\bottomrule
\endlastfoot
均值基线 & NA & 1.0020 & 0.0463 \\
均值基线 & PA & 1.0002 & 0.0836 \\
均值基线 & LS & 1.0001 & 0.0605 \\
原共享 RF & NA & 1.0027 & 0.0473 \\
原共享 RF & PA & 0.9986 & 0.0857 \\
原共享 RF & LS & 0.9991 & 0.0578 \\
所选程序 & NA & 1.0044 & 0.0481 \\
所选程序 & PA & 0.9997 & 0.0822 \\
所选程序 & LS & 0.9976 & 0.0572 \\
\end{longtable}
\tabnote{NA 为消极情绪，PA 为积极情绪，LS 为生活满意度。五个外层折，每个外层折内部重新执行五折候选选择。宏平均：基线 1.0007，原 RF 1.0001，所选程序 1.0006。这些是归一化误差，不是存储单位下的留出 RMSE。由于此前开发和留出结果已知，整个扩展均属探索性。标准差仅作描述。}
\endgroup
